\section{Cartan models and the canonical height}
\label{sec:models}

Throughout this section, $G$ is a connected, simply connected,
real-triangulable, unimodular Lie group, with Lie algebra $\mathfrak g$.
We fix left invariant Riemannian metrics on every group under consideration.
The notation $|x|_J$ denotes distance from the identity in the intrinsic
metric of $J$. We construct nilpotent groups $N,D$ whose multiplication
maps onto $G$, identify their height coordinates with the canonical quotient
of Definition~\ref{def:canonical-height}, and establish the polynomial
coordinate estimates needed to retain the nonsemisimple parts of the
adjoint action.

The fixed metrics are complete. Indeed, choose a relatively compact identity
neighborhood containing a metric ball of positive radius. A Cauchy
sequence eventually lies in one left translate of that ball, hence in
a compact set; a convergent subsequence then forces the whole sequence
to converge. By the Hopf--Rinow theorem, closed bounded sets are compact
and any two points can be joined by a minimizing geodesic.

\subsection{The Cartan decomposition and its multiplication map}

A Cartan subalgebra of a solvable Lie algebra is a nilpotent,
self-normalizing subalgebra.  Choose one, denoted by $\mathfrak d$.
For a real character $\gamma:\mathfrak d\to\R$, its generalized weight
space in $\mathfrak g$ is
\[
 \mathfrak g_\gamma
 =\{u\in\mathfrak g:
   (\ad(v)-\gamma(v)\id)^{\dim\mathfrak g}u=0
   \text{ for every }v\in\mathfrak d\}.
\]
Only characters with a nonzero space are called weights.  The exponent
in this definition can be replaced by any sufficiently large fixed integer.
In particular, this definition allows a nilpotent part on each weight
space and does not require $\mathfrak d$ to be abelian.

\begin{lemma}[Cartan model]\label{lem:cartan-model}
There is a real generalized weight decomposition
\begin{equation}\label{eq:cartan-weight-decomposition}
 \mathfrak g=\bigoplus_\gamma\mathfrak g_\gamma,
 \qquad \mathfrak g_0=\mathfrak d,
 \qquad
 [\mathfrak g_\gamma,\mathfrak g_\eta]
       \subseteq\mathfrak g_{\gamma+\eta}.
\end{equation}
Let $\mathfrak n$ be the subalgebra generated by the spaces with nonzero
weight.  Then $\mathfrak n$ is a nilpotent ideal contained in
$[\mathfrak g,\mathfrak g]$, and $\mathfrak g=\mathfrak n+\mathfrak d$.
There are homogeneous generators $X_1,\ldots,X_\nu$ of $\mathfrak n$ whose
images form a basis of $\mathfrak n/[\mathfrak n,\mathfrak n]$ and whose
weights $\alpha_1,\ldots,\alpha_\nu$ are all nonzero.

The groups $N=\exp_G(\mathfrak n)$ and $D=\exp_G(\mathfrak d)$ are closed
and simply connected, and $N$ is normal in $G$.  With the conjugation
action of $D$ on $N$, multiplication defines a surjective Lie homomorphism
and submersion
\begin{equation}\label{eq:cartan-multiplication}
 q:N\rtimes D\longrightarrow G,\qquad q(n,d)=nd.
\end{equation}
Its kernel is
\begin{equation}\label{eq:cartan-overlap-kernel}
 \ker q=\{(a^{-1},a):a\in N\cap D\}.
\end{equation}
Moreover, $N\cap D=\exp_G(\mathfrak n\cap\mathfrak d)$.
Every compact subset of $G$ has lifts in a compact subset of $N\rtimes D$.
Finally, $N=1$ if and only if $G$ is nilpotent.
\end{lemma}

\begin{proof}
We give the decomposition argument to specify precisely which parts of
the Cartan action may be nonsemisimple; see also
\cite[Section~4.D]{CornulierTessera2017} for the Cartan grading.
For $v,u\in\mathfrak d$, nilpotence of $\mathfrak d$ gives
\[
 (\ad_{\End(\mathfrak g)}(\ad v))^b(\ad u)
   =\ad((\ad v)^b u)=0
\]
for sufficiently large $b$.  Consequently $\ad u$ preserves every
generalized eigenspace of $\ad v$.  Indeed, between two primary blocks
with distinct eigenvalues, the commutator with $\ad v$ is invertible,
so an operator killed by a power of that commutator has no such block.
All eigenvalues are real by real triangulability.  Successively taking
the primary decompositions for a basis of $\mathfrak d$ therefore gives
simultaneous invariant primary spaces.  On each such space all diagonal
entries of a triangular representation agree on that basis, and hence
agree as linear forms on $\mathfrak d$.  Their common form $\gamma$ is a
character, since diagonals of commutators vanish.  This proves the direct
sum decomposition into the spaces in
\eqref{eq:cartan-weight-decomposition}.

The derivation identity shows that weights add under brackets: applying
$\ad v-(\gamma(v)+\eta(v))\id$ repeatedly to $[x,y]$, with
$x\in\mathfrak g_\gamma$ and $y\in\mathfrak g_\eta$, gives a binomial
sum involving powers of $\ad v-\gamma(v)\id$ on $x$ and of
$\ad v-\eta(v)\id$ on $y$.  Every summand vanishes at a sufficiently
large power.  Nilpotence of $\mathfrak d$ gives
$\mathfrak d\subseteq\mathfrak g_0$.  If this inclusion were strict,
the induced action of $\mathfrak d$ on $\mathfrak g_0/\mathfrak d$
would consist of nilpotent operators.  Engel's theorem would give a
nonzero fixed class.  Any representative of that class normalizes
$\mathfrak d$, contrary to self-normalization.  Thus
$\mathfrak g_0=\mathfrak d$.

The nonzero weight spaces are $\mathfrak d$-invariant.  Their generated
subalgebra $\mathfrak n$ is therefore $\mathfrak d$-invariant, and
$\mathfrak g=\mathfrak n+\mathfrak d$ makes it an ideal.  If
$\gamma\ne0$, choose $v\in\mathfrak d$ with $\gamma(v)\ne0$.
The restriction of $\ad v$ to $\mathfrak g_\gamma$ is invertible, so
$\mathfrak g_\gamma\subseteq[\mathfrak g,\mathfrak g]$.
In a simultaneous real triangularization of $\ad(\mathfrak g)$,
every element of $\ad([\mathfrak g,\mathfrak g])$ is strictly upper
triangular.  Engel's theorem applied to the derived algebra proves
its nilpotence, and hence that of $\mathfrak n$.

The images of the generating nonzero weight spaces span
$\mathfrak n/[\mathfrak n,\mathfrak n]$, and its induced weight
decomposition has no zero summand.  Choose a homogeneous basis of this
quotient and homogeneous lifts $X_i$.  These lifts generate
$\mathfrak n$: if $\mathfrak h$ is their generated subalgebra and
$\mathfrak n^j$ denotes the lower central series, then
$\mathfrak n=\mathfrak h+\mathfrak n^2$.  Induction, using brackets,
gives $\mathfrak n=\mathfrak h+\mathfrak n^j$ for every $j$, and
nilpotence concludes the argument.

We now pass from Lie algebras to groups. Dixmier's real-root criterion implies that
$\exp_G:\mathfrak g\to G$ is an analytic diffeomorphism
\cite[Theorem~3]{Dixmier1957}.  The inclusion of either nilpotent
subalgebra integrates from its simply connected Baker--Campbell--Hausdorff
group, and this homomorphism is the restriction of $\exp_G$ in
exponential coordinates.  It is injective, with closed image, because
$\exp_G$ is a diffeomorphism and each subalgebra is a closed vector
subspace.  Injectivity of $\exp_G$ also gives the stated description
of $N\cap D$.  This proves the claims about $N,D$, including normality
of $N$.

The group law in the domain of $q$ is
\[
 (n,d)(n',d')=(n(dn'd^{-1}),dd'),
\]
which verifies the homomorphism assertion.  Its derivative at the
identity is $(x,v)\mapsto x+v$, which is onto.  Its image is therefore
an open subgroup of the connected group $G$, so it is all of $G$;
translation also makes $q$ a submersion everywhere.  The kernel formula
follows directly from $nd=1$.  Local smooth sections of a submersion,
restricted to relatively compact neighborhoods, give compact sets of
lifts; a finite cover proves the claim for an arbitrary compact set.

If $\mathfrak n=0$, then $\mathfrak g=\mathfrak d$ is nilpotent.
Conversely, in a nilpotent $\mathfrak g$ all adjoint operators are
nilpotent.  The decomposition just proved then has only weight zero,
so $\mathfrak g=\mathfrak d$ and $\mathfrak n=0$.
\end{proof}

The map $q$ will let us estimate paths using a nilpotent $D$ coordinate.
Its kernel in \eqref{eq:cartan-overlap-kernel} is retained throughout;
the argument does not require a complementary subgroup to $N$.
To control conjugation by a $D$ element of intrinsic length $O(r)$,
we need polynomial bounds on its logarithmic coordinates. We establish
those bounds next, together with the rescaling of $D$ used for large charts.

\subsection{Intrinsic coordinates and the graded Cartan limit}

Write $\mathfrak d^1=\mathfrak d$ and
$\mathfrak d^{j+1}=[\mathfrak d,\mathfrak d^j]$ for the lower central
series.  Choose vector subspaces $V_j$ such that
\begin{equation}\label{eq:d-filtration-splitting}
 \mathfrak d=V_1\oplus\cdots\oplus V_s,
 \qquad \mathfrak d^j=V_j\oplus\cdots\oplus V_s.
\end{equation}
Thus $V_1$ represents $V=\mathfrak d/[\mathfrak d,\mathfrak d]$.
Fix Euclidean norms on these spaces.  Define
$\delta_r(\sum_j v_j)=\sum_j r^jv_j$ for $r>0$.  The associated
graded bracket on the same vector space is given, for $v_i\in V_i$
and $v_j\in V_j$, by
\[
 [v_i,v_j]_\infty=\operatorname{pr}_{i+j}[v_i,v_j].
\]
Here $\operatorname{pr}_k=0$ for $k>s$.
Let $D_\infty$ be its simply connected nilpotent group and set
\begin{equation}\label{eq:d-delta-definition}
 \Delta_r(\exp_\infty v)=\exp_D(\delta_r v).
\end{equation}
When $D=1$, take $s=1$ and $V_1=0$; the assertions below are immediate.

These weighted coordinates and graded limits belong to the classical
large-scale geometry of nilpotent groups; see
\cite{Guivarch1973,Pansu1983}. The maps $\Delta_r$ will let us place
growing $D$-charts on the fixed domain $D_\infty$. We give the finite
Baker--Campbell--Hausdorff estimates needed here.

\begin{lemma}[Nilpotent coordinates and scaling]\label{lem:d-scaling}
There is a constant $C\ge1$, depending only on $D$, the metric, and the
splitting, such that for $d=\exp_D(\sum_jv_j)$,
\begin{align}
 \|v_j\|&\le C(1+|d|_D)^j \quad(1\le j\le s),
       \label{eq:d-coordinate-upper}\\
 |d|_D&\le C\left(1+\max_{1\le j\le s}\|v_j\|^{1/j}\right).
       \label{eq:d-coordinate-reverse}
\end{align}
The group laws transported from $D$ by $\Delta_r$ converge, with all
derivatives on compact coordinate sets, to the law of $D_\infty$ as
$r\to\infty$.  Every compact $K\subset D_\infty$ satisfies
$\sup_{z\in K}|\Delta_r z|_D=O_K(r)$ for $r\ge1$.
Conversely, for every fixed $C_0$, all
$\Delta_r^{-1}\{d:|d|_D\le C_0r\}$, $r\ge1$, lie in one compact
coordinate set.
If $z_r,z'_r$ belong to a fixed compact coordinate set and their coordinate
difference tends to zero, then
\begin{equation}\label{eq:d-scaling-equicontinuity}
 \dist_D(\Delta_r z_r,\Delta_r z'_r)=o(r).
\end{equation}
For $v\in V_1$ one also has
\begin{equation}\label{eq:d-horizontal-flow-limit}
 \Delta_r^{-1}\bigl(\Delta_r z\exp_D(rt v)\bigr)
     \longrightarrow z\exp_\infty(tv),
\end{equation}
uniformly when $z,t,v$ range over fixed compact sets.
Haar measures on $D$ and $D_\infty$ are Lebesgue measure in these
exponential coordinates, up to fixed scalars.  If both scalars are
chosen equal and $Q=\sum_jj\dim V_j$, then $\Delta_r$ multiplies Haar
volume by $r^Q$, and
\begin{equation}\label{eq:d-polynomial-volume}
 c t^Q\le \vol_D\{d:|d|_D\le t\}\le C t^Q \qquad(t\ge1)
\end{equation}
for some $c>0$.  In particular $D$ has doubling Haar measure:
\begin{equation}\label{eq:nilpotent-doubling}
 \vol_D B_D(1,2t)\le C\vol_D B_D(1,t)\qquad(t>0).
\end{equation}
Every conclusion of this lemma holds for an arbitrary connected,
simply connected nilpotent Lie group $L$ equipped with a left invariant
Riemannian metric and a splitting of its lower central filtration,
with $L$ in place of $D$ and constants allowed to depend on $L$.
\end{lemma}

\begin{proof}
We first bound the coordinates of a short path. For the reverse estimate,
we construct uniformly bounded products of degree-one exponentials in
the rescaled groups. The remaining assertions follow from these
coordinates and the Haar Jacobian.

A path of length at most $t$ splits its endpoint into a product of
at most $C(1+t)$ increments from a fixed compact neighborhood.  Their
logarithms lie in a fixed bounded set.  In a product of $a$ exponentials,
the degree-$b$ part of its formal logarithm has a sum of absolute
coefficients at most $C_b(1+a)^b$.  To see the bound, expand the product
of formal exponentials through degree $s$ in the noncommuting variables,
and then use the finite expansion of $\log(1+u)$ in those degrees.
The degree-$b$ coefficient sum of the product is at most $a^b/b!$;
the finitely many compositions of $b$ appearing in the logarithm give
the asserted bound.  Applying the degree-$b$ Dynkin projection, which
sends an associative monomial to its iterated bracket and acts by
multiplication by $b$ on homogeneous Lie polynomials, converts this
expansion to brackets with only a degree-dependent coefficient bound.
A bracket involving more than $j$ increments belongs to
$\mathfrak d^{j+1}$ and has zero $V_j$ projection.  Evaluating the
remaining finitely many types of brackets on the bounded logarithms
therefore bounds this projection by $C_j(1+a)^j$.  This proves
\eqref{eq:d-coordinate-upper}.

For the converse, we first record the limit of the transported laws.
The inclusion
$[V_i,V_j]\subseteq\bigoplus_{k\ge i+j}V_k$ shows that the coefficient
of a bracket from $V_i\times V_j$ into $V_k$, when transported by
$\delta_r$, is multiplied by $r^{i+j-k}$.  For $r\ge1$ these
coefficients are bounded and converge to the associated graded bracket.
The finite Baker--Campbell--Hausdorff formula proves convergence of the
group laws and all their derivatives on compact sets.  Denote the
transported group by $D_r$; its underlying exponential coordinate space
is the fixed vector space $\bigoplus_jV_j$.

We now produce uniformly bounded products of degree-one exponentials
covering a fixed coordinate neighborhood.  Choose a basis
$e_1,\ldots,e_a$ of $V_1$.  These vectors generate the graded Lie algebra:
by the definition of the lower central series, the quotient
$\mathfrak d^j/\mathfrak d^{j+1}$ is spanned by brackets of length $j$
of classes in $\mathfrak d/\mathfrak d^2$.
In $D_\infty$, let $H$ consist of all finite products of
$\exp_\infty(t e_i)$ and put
\[
 L=\Span\{\Ad_h e_i:h\in H,\ 1\le i\le a\}.
\]
This space contains the $e_i$ and is invariant under every
$\Ad_{\exp_\infty(t e_i)}$.  Differentiation makes it invariant under
$\ad e_i$, so it contains all their brackets and equals the entire
graded Lie algebra.  Choose $b=\dim D$ such conjugated vectors
$\Ad_{h_\ell}e_{i_\ell}$ forming a basis, where each $h_\ell$ is a
fixed finite product of degree-one exponentials.  The map
\[
 \Phi_\infty(t_1,\ldots,t_b)
   =\prod_{\ell=1}^b
      h_\ell\exp_\infty(t_\ell e_{i_\ell})h_\ell^{-1}
\]
is the identity at $t=0$ and has these basis vectors as its derivative
columns there.

Use the same fixed words for $h_\ell$ in $D_r$ to define
$\Phi_r$.  The maps $\Phi_r$ converge with their derivatives to
$\Phi_\infty$ on compact parameter sets.  The inverse function theorem
therefore supplies a fixed parameter ball and a fixed coordinate
neighborhood $U$ of the identity contained in their images for all
sufficiently large $r$.  More explicitly, choose a parameter ball on
which the derivative of $\Phi_\infty$ differs by less than one quarter of the
smallest singular value of $D\Phi_\infty(0)$ from that derivative;
the corresponding estimate with one half holds for $\Phi_r$.
The usual contraction proof of the inverse function theorem then gives
one common image ball.  Every element of $U$ is thus a product of a
fixed number of $\exp_r(c e_i)$, with all $|c|$ bounded by one constant.
The conjugators here are themselves fixed words in these exponentials.

A fixed number of powers of $U$ covers the coordinate box
$\{\sum_jv_j:\|v_j\|\le1\}$, uniformly in $r$:
choose an integer $M$ large enough that every vector $v$ in this box
satisfies $v/M\in U$, and use
$\exp_r(v)=(\exp_r(v/M))^M$.
Under $\Delta_r$, each degree-one factor becomes $\exp_D(rc e_i)$,
whose intrinsic length is at most $r|c|\|e_i\|_D$, where the last norm
is the Riemannian norm at the identity.  The resulting
products have length at most $Cr$.  Enlarging $C$ handles the remaining
compact interval of $r\ge1$.  Given arbitrary $v_j$, apply this bound
with $r=\max(1,\max_j\|v_j\|^{1/j})$ to obtain
\eqref{eq:d-coordinate-reverse}.

For the equicontinuity assertion, the logarithm of the difference of
$z_r$ and $z'_r$ in the transported law is a vector
$\xi_r=\sum_j\xi_{r,j}$ tending to zero.  The logarithm of the
corresponding difference in $D$ is $\delta_r\xi_r$.  Thus
\eqref{eq:d-coordinate-reverse} gives
\[
 \frac{\dist_D(\Delta_r z_r,\Delta_r z'_r)}r
   \le \frac Cr+C\max_j\|\xi_{r,j}\|^{1/j}\longrightarrow0.
\]
The compact length bound follows from the same coordinate estimate.
The converse compactness assertion follows from
\eqref{eq:d-coordinate-upper} after dividing the $V_j$ coordinate by
$r^j$.
Equation~\eqref{eq:d-horizontal-flow-limit} follows directly from
convergence of the transported product, since $\delta_r(tv)=rtv$ on
$V_1$.

Finally, the differential of the exponential map, after left translation
to the identity, is a polynomial in $\ad v$ with constant term $\id$.
Since $\ad v$ is nilpotent, this polynomial has determinant one.  This proves the
claims about Haar measure on both nilpotent groups.  The linear
coordinate map $\delta_r$ has determinant $r^Q$.  The coordinate bounds
\eqref{eq:d-coordinate-upper} place a radius-$t$ ball in a product of
Euclidean balls with radii $C(1+t)^j$, proving the upper volume bound.
Conversely, \eqref{eq:d-coordinate-reverse} puts the coordinate box
$\{\sum_jv_j:\|v_j\|\le(t/(2C))^j\}$ inside the radius-$t$ ball
when $t\ge2C$.  This box has volume a fixed positive constant times
$t^Q$.  Decreasing that constant handles $1\le t\le2C$, and proves
\eqref{eq:d-polynomial-volume}.  These two bounds give doubling for
$t\ge1$.  At sufficiently small radii a smooth coordinate chart and
the Riemannian metric give volume comparable to $t^{\dim D}$; the
remaining compact interval of radii is handled by monotonicity and
positivity of ball volume.  This proves
\eqref{eq:nilpotent-doubling}.  The proof used only simple connectivity
and nilpotence, so it applies without change to the group $L$ in the
last assertion.
\end{proof}

\subsection{Canonical height and the adjoint action}

We now identify the canonical quotient $W$ in terms of the Cartan weights.
This also distinguishes the exact scalar action of height from actual
conjugation by $D$: the latter may have a unipotent factor, which the
preceding coordinate estimate will bound polynomially in intrinsic length.
Choose a simultaneous real
triangularization of $\ad(\mathfrak g)$, and let $\chi_1,\ldots,\chi_b$
be its diagonal characters, including repetitions and zero characters.
Thus
\[
 \mathfrak k=\bigcap_{j=1}^b\ker\chi_j,\qquad W=\mathfrak g/\mathfrak k.
\]
These are the adjoint characters, rather than all real homomorphisms
on $\mathfrak g$.  Their multiset is independent of the triangularization,
as the multiset of composition factors of the adjoint representation.
Every $\chi_j$ vanishes on $[\mathfrak g,\mathfrak g]$, so the quotient
is abelian.  Write $\pi_*:\mathfrak g\to W$ for the quotient and
$\pi:G\to W$ for its integrated homomorphism.  As usual, we also write
$\pi$ for induced linear maps when their domain is clear.

\begin{lemma}[Canonical height and volume-preserving actions]
\label{lem:canonical-height}
Each generator weight $\alpha_i$ vanishes on
$[\mathfrak d,\mathfrak d]+(\mathfrak d\cap\mathfrak n)$, and inclusion
of $\mathfrak d$ into $\mathfrak g$ induces an isomorphism
\begin{equation}\label{eq:canonical-cartan-height}
 W\cong V\Big/\bigcap_{i=1}^\nu\ker\alpha_i,
 \qquad V=\mathfrak d/[\mathfrak d,\mathfrak d].
\end{equation}
Under this identification, $\pi$ kills $N$ and its restriction to $D$
is the displayed quotient of its abelianization.  When $\nu=0$, the
intersection on the right is all of $V$, and $W=0$; this occurs exactly
when $G$ is nilpotent.

Write $\mathfrak n_\gamma=\mathfrak n\cap\mathfrak g_\gamma$ and view
every weight occurring in $\mathfrak n$ as a linear form on $W$.
For fixed norms, there are constants $C,M$ such that, for every $d\in D$
and every such weight $\gamma$,
\begin{equation}\label{eq:d-polynomial-unipotent}
 \Ad_d|_{\mathfrak n_\gamma}
     =e^{\gamma(\pi d)}U_\gamma(d),\qquad
 \|U_\gamma(d)\|+\|U_\gamma(d)^{-1}\|
     \le C(1+|d|_D)^M,
\end{equation}
where $U_\gamma(d)$ is unipotent.  For $w\in W$, the linear map
\begin{equation}\label{eq:exact-weight-dilation}
 S_w|_{\mathfrak n_\gamma}=e^{\gamma(w)}\id
\end{equation}
is a Lie algebra automorphism and hence a group automorphism of $N$.
Both $S_w$ and conjugation by $d$ preserve Haar measure on $N$.
Consequently product Haar measure $dn\,dd$ is Haar measure on $N\rtimes D$.

The height map on $D_\infty$ is the homomorphism induced by
$V_1\cong V\to W$, and it satisfies
\begin{equation}\label{eq:d-height-scaling}
 \pi(\Delta_r z)=r\pi(z)
 \qquad(z\in D_\infty, r>0).
\end{equation}
\end{lemma}

\begin{proof}
The restrictions of the diagonal characters $\chi_j$ to $\mathfrak d$
are exactly the generalized weights of its action on $\mathfrak g$,
with their multiplicities.  All $\chi_j$ vanish on $\mathfrak n$ by
Lemma~\ref{lem:cartan-model}.  In particular every Cartan weight
vanishes on $\mathfrak d\cap\mathfrak n$, and, being a character, also
on $[\mathfrak d,\mathfrak d]$.

Homogeneous bracket words in the $X_i$ span $\mathfrak n$.
Their weights are sums of the $\alpha_i$, whereas every $\alpha_i$
itself occurs.  The weights on $\mathfrak d=\mathfrak g_0$ are zero.
It follows that
\[
 \mathfrak k\cap\mathfrak d=\bigcap_i\ker\alpha_i
       \quad\text{as subspaces of }\mathfrak d.
\]
Since $\mathfrak n\subseteq\mathfrak k$ and
$\mathfrak g=\mathfrak n+\mathfrak d$, inclusion gives
$\mathfrak g/\mathfrak k\cong\mathfrak d/(\mathfrak k\cap\mathfrak d)$, and passing
through $V$ proves \eqref{eq:canonical-cartan-height}.
It also shows that every weight of $\mathfrak n$ descends to $W$.
The statements about the integrated maps follow from simple connectivity.
If $\nu=0$, Lemma~\ref{lem:cartan-model} gives nilpotence and $W=0$.
Conversely, if $W=0$, all adjoint characters vanish, so every adjoint
operator is strictly triangular and Engel's theorem makes $G$ nilpotent.

For $d=\exp_D(v)$ and $b_\gamma=\dim\mathfrak n_\gamma$, write
\[
 \ad(v)|_{\mathfrak n_\gamma}
    =\gamma(\pi_*v)\id+T_\gamma(v),
 \qquad T_\gamma(v)^{b_\gamma}=0.
\]
The map $v\mapsto T_\gamma(v)$ is linear, and exponentiation gives
\[
 U_\gamma(d)=
   \sum_{j=0}^{b_\gamma-1}\frac{T_\gamma(v)^j}{j!}.
\]
Replacing $v$ by $-v$ gives its inverse.  By
\eqref{eq:d-coordinate-upper}, $\|v\|$ is bounded by a polynomial in
$1+|d|_D$.  The two finite sums therefore prove
\eqref{eq:d-polynomial-unipotent}, uniformly over the finitely many
weight spaces, including the zero weight when it occurs.

For $v\in\mathfrak d$, the action of $\ad v$ on
$\mathfrak g/\mathfrak n\cong
\mathfrak d/(\mathfrak d\cap\mathfrak n)$ is nilpotent.  Unimodularity
of $G$ consequently gives
\begin{equation}\label{eq:trace-on-n-zero}
 \tr(\ad v|_{\mathfrak n})
 =\tr(\ad v|_{\mathfrak g})
   -\tr(\ad v|_{\mathfrak g/\mathfrak n})=0.
\end{equation}
Thus $\det(\Ad_d|_{\mathfrak n})=1$.  Bracket addition in
\eqref{eq:cartan-weight-decomposition} proves that $S_w$ preserves
brackets.  Choosing $v\in\mathfrak d$ over $w$ in
\eqref{eq:trace-on-n-zero} also gives
\[
 \det S_w
 =\exp\left(\sum_\gamma
                  (\dim\mathfrak n_\gamma)\gamma(w)\right)=1.
\]
An automorphism of a simply connected nilpotent group changes Haar
measure by its Lie algebra determinant, so both asserted actions
preserve Haar measure.  As $N$ and $D$ are unimodular and the action
preserves Haar measure on $N$, direct substitution in the semidirect
product law proves that $dn\,dd$ is left Haar measure on $N\rtimes D$.

Finally, $\pi_*$ annihilates $\mathfrak d^2$, so in the splitting
\eqref{eq:d-filtration-splitting} it depends only on $V_1$.
The same projection is a homomorphism on the associated graded group.
Equation~\eqref{eq:d-height-scaling} now follows from
$\delta_r|_{V_1}=r\id$.
\end{proof}

The construction supplies the canonical height using any chosen Cartan
subalgebra.  It also separates the exponential factors
$e^{\gamma(\pi d)}$ from polynomial factors in the intrinsic $D$ length.
The weighted boxes developed next will keep these polynomial factors
while recording exponential volume precisely.
